\documentclass[conference]{IEEEtran}
% IEEE conference paper template (IEEEtran v1.8b)

\usepackage{cite}
\usepackage{amsmath,amssymb,amsfonts}
\usepackage{graphicx}
\usepackage{textcomp}
\usepackage{booktabs}
\usepackage{multirow}
\usepackage{array}
\usepackage{url}
\usepackage{hyperref}
\hypersetup{colorlinks=true,linkcolor=blue,citecolor=blue,urlcolor=blue}

\newcommand{\keywords}[1]{\par\addvspace\medskipamount\noindent\textbf{Keywords:} #1}

\begin{document}

\title{Performance Evaluation of the RSA Cryptosystem in\\Web Applications: Standard vs.\ Optimized Implementations}

\author{
\IEEEauthorblockN{Bhavesh Patil, Devendra Patil, Sakti vala}
\IEEEauthorblockA{Department of Information Technology\\
C. K. Pithawala College of Engineering\\
Surat, Gujarat, India\\
}
}

\maketitle

\begin{abstract}
Due to rapid technological growth, online applications now face serious security challenges in data transmission, reception, and user management. The Rivest--Shamir--Adleman (RSA) cryptosystem is still widely used for secure key exchange, digital signatures, and identity verification. However, traditional RSA is slow, consumes high computational power, and struggles with efficient key management in web environments. This study analyzes RSA's performance and security in web applications and highlights its weaknesses against complex mathematical and side-channel attacks. To improve efficiency, we explore optimization techniques such as the Chinese Remainder Theorem (CRT) for decryption, sliding-window modular exponentiation, Montgomery multiplication, key caching, and hybrid RSA+AES-256-GCM envelope encryption. Experimental results indicate that key caching provides the largest practical speedup (43--93$\times$), CRT-based decryption yields a measurable 1.23$\times$ improvement at 2048-bit keys, and hybrid encryption is 8.3--398$\times$ faster than pure RSA for bulk data. The study finds that hybrid RSA frameworks are necessary for modern web security, especially as quantum computing continues to develop and create new threats to data protection.
\end{abstract}

\keywords{RSA Algorithm, Web Application Security, Performance Optimization, Hybrid End-to-End Encryption, Chinese Remainder Theorem, Post-Quantum Cryptography}

\section{Introduction}
\label{sec:intro}

Because information technology is growing so quickly, web apps are now a part of everyday life and work. This phenomenon has led to a demand for strong data protection in areas like e-commerce, social networks, and online finance. However, this growth also poses significant cybersecurity risks, particularly in terms of safeguarding users' privacy during data transmission and processing on the cloud. HTTPS and other classic transport layer encryption methods have been crucial for keeping data safe, but they often leave server-side data processing and storage open to attack. As a result, creating effective and reliable encryption systems has become a key focus for researchers to address the high needs of today's web systems.

The Rivest--Shamir--Adleman (RSA) cryptosystem~\cite{rivest1978} is a fundamental asymmetric encryption technique that protects information by making it difficult to factor massive numbers. It is widely used to keep personal information safe by allowing secure key exchange, digital signatures, and identity verification. Even though many people use it, typical RSA encryption takes a lot of computing power, which makes processing slow and key management challenging. Due to these speed issues, RSA is not a viable way to encrypt large amounts of data directly. Instead, it is mostly used to protect key exchanges.

The RSA algorithm has more than just computational limits; it also has significant and evolving security problems. It is subject to attacks like modulus decomposition, co-modulus attacks, and small exponent assaults because of the way its internal math works~\cite{boneh1999}. Advanced cryptanalysis, especially lattice-based approaches developed by Coppersmith~\cite{coppersmith1997}, has shown that RSA variants can be attacked in polynomial time under certain conditions. What's even more concerning is that quantum computing could end RSA; Shor's algorithm~\cite{shor1997} can solve the integer factorization problem in polynomial time.

Modern systems must use more than one encryption method because a single method generally doesn't meet both the high performance and strong security needs of sophisticated online applications. The hybrid end-to-end encryption (E2EE) architecture effectively balances the rapid processing of symmetric algorithms such as AES with the robust key management of asymmetric algorithms.

This study offers an exhaustive assessment of the RSA algorithm's performance and security in contemporary online applications. We examine algorithm-level optimizations, including CRT-based decryption, sliding-window exponentiation, Montgomery multiplication~\cite{montgomery1985}, and key caching, alongside dynamic load-aware resource allocation designed to expedite encryption during peak user activity. We also examine the application-layer payload protection of RSA in hybrid E2EE architectures. Finally, the paper discusses the proactive shift to Post-Quantum Cryptography (PQC)~\cite{nist2024} to ensure that web application security remains future-proof.

\section{Literature Review}
\label{sec:litreview}

\subsection{Improving RSA Performance and Solving Speed Issues}

Although the RSA algorithm~\cite{rivest1978} is essential to digital security, its reliance on large integer factorization makes it computationally demanding. When it comes to large-scale data transmission and web applications with high traffic, RSA frequently degrades system performance. Recent research focuses on improving resource management and algorithm-level optimization to address this. Through the use of fast modular exponentiation methods and optimal parameter selection, the SMM-RSA technique and similar approaches enhance performance. These enhancements can lower memory usage and boost encryption efficiency by about 30\% when paired with dynamic, load-aware resource allocation. Using big data platforms like Hadoop for distributed RSA processing is another way to handle large-scale encryption tasks more effectively.

\subsection{Classical Cryptanalysis and Lattice-Based Attacks}

While performance optimization is important, RSA's resistance to advanced cryptanalysis remains a major research area. Historically, RSA has faced attacks targeting its mathematical structure, such as Wiener's attack on small private exponents~\cite{wiener1990} and partial key exposure attacks. A major breakthrough in RSA cryptanalysis came from Coppersmith's lattice-based method~\cite{coppersmith1997}, which transforms the RSA-breaking problem into finding short vectors in a lattice. This approach is effective when partial key information is known. Boneh, DeMillo, and Lipton~\cite{boneh2001} demonstrated that even computational faults during RSA operations can leak key information, motivating the inclusion of fault-attack protection in our optimized implementation.

\subsection{Hybrid End-to-End Encryption Architectures}

Modern research supports hybrid E2EE models because no single encryption method can provide both high speed and strong security. Hybrid systems combine symmetric encryption algorithms like AES-GCM, which are fast for encrypting large data, with asymmetric algorithms like RSA or ECC for secure key exchange and authentication~\cite{menezes2001}. This method protects data at both the application and network levels and addresses limitations of traditional transport protocols like HTTPS.

\subsection{Quantum Computing Threat and Post-Quantum Cryptography}

Quantum computing is considered the most serious future threat to RSA. Shor's algorithm~\cite{shor1997} proves that a sufficiently powerful quantum computer could break RSA-3072 in polynomial time. Because of this, there is global agreement on transitioning to Post-Quantum Cryptography (PQC). Lattice-based schemes such as CRYSTALS-Kyber~\cite{kyber2023} and CRYSTALS-Dilithium~\cite{dilithium2022} rely on the Learning With Errors problem, which is resistant to both classical and quantum attacks.

\section{Methodology}
\label{sec:methodology}

\subsection{Research Design}

This study follows a quantitative experimental approach. Two independent RSA implementations were developed from scratch in Python: a standard (unoptimized) baseline and an optimized variant incorporating multiple algorithmic improvements. Both implementations share the same mathematical foundation but differ in the computational strategies used for modular exponentiation, key management, and decryption. A suite of six benchmark tests was then designed and executed against both implementations to measure correctness, speed, memory consumption, scalability, hybrid encryption efficiency, and concurrent throughput.

\subsection{Implementation Architecture}

The project is organized into four modules totaling approximately 1,690 lines of Python code:

\begin{itemize}
    \item \textbf{standard\_rsa/}: A baseline RSA implementation with no algorithmic optimizations. This module implements key generation using the Miller-Rabin primality test, textbook RSA encryption and decryption using $c = m^e \bmod n$ and $m = c^d \bmod n$, naive square-and-multiply modular exponentiation, and PKCS\#1 v1.5-style padding~\cite{pkcs1_2016}. Every function call generates fresh keys with no caching.

    \item \textbf{optimized\_rsa/}: An enhanced RSA implementation incorporating eight optimizations: (1)~CRT for decryption~\cite{garner1959}, splitting full-size exponentiation into two half-size operations; (2)~sliding-window modular exponentiation with a 4-bit window; (3)~Montgomery multiplication~\cite{montgomery1985}; (4)~key caching via a \texttt{KeyCache} class; (5)~batch encrypt/decrypt functions; (6)~pre-computed CRT parameters ($d_p$, $d_q$, $q_{\mathrm{inv}}$); (7)~hybrid RSA+AES-256-GCM encryption; and (8)~fault-attack protection~\cite{boneh2001}.

    \item \textbf{tests/}: Six comprehensive test files evaluating both implementations across multiple performance dimensions.

    \item \textbf{run\_all\_tests.py}: A unified test runner executing all six test suites sequentially.
\end{itemize}

\subsection{Mathematical Foundations of the Optimizations}

The standard RSA decryption computes $m = c^d \bmod n$ using a single full-size modular exponentiation where $d$ is roughly the same bit-length as $n$. For a 2048-bit key, this means a 2048-bit exponent is processed one bit at a time.

The CRT optimization reformulates this computation. Given prime factors $p$ and $q$ of $n$, the private exponent $d$ is decomposed into two half-size exponents:
\begin{align}
    d_p &= d \bmod (p - 1) \label{eq:dp} \\
    d_q &= d \bmod (q - 1) \label{eq:dq}
\end{align}
Decryption proceeds as two independent half-size modular exponentiations followed by Garner's recombination~\cite{garner1959}:
\begin{align}
    m_1 &= c^{d_p} \bmod p \label{eq:m1} \\
    m_2 &= c^{d_q} \bmod q \label{eq:m2} \\
    h &= q_{\mathrm{inv}} \cdot (m_1 - m_2) \bmod p \label{eq:h} \\
    m &= m_2 + h \cdot q \label{eq:m}
\end{align}

Since modular exponentiation complexity is $O(n^2 \log n)$ for $n$-bit operands using standard multiplication, halving the operand size reduces each exponentiation by approximately 4$\times$, yielding a combined speedup of roughly 3--4$\times$.

The sliding-window method pre-computes a table of odd powers $g^1, g^3, g^5, \ldots, g^{2^w - 1}$ for window size $w$, then scans the exponent in windows of up to $w$ bits. For a 4-bit window, this reduces the average number of multiplications from $n$ to approximately $n / (w + 1) + 2^{w-1} - 1$.

Montgomery multiplication~\cite{montgomery1985} replaces expensive modular reduction (division by $n$) with shift-based reduction using auxiliary modulus $R = 2^k$. Numbers are converted to Montgomery form ($\bar{x} = x \cdot R \bmod n$), and the REDC algorithm replaces division by $n$ with division by $R$ (a simple right shift).

\subsection{Hybrid Encryption Architecture}

The hybrid RSA+AES-GCM architecture follows the envelope encryption model. For each encryption operation:
\begin{enumerate}
    \item A random 256-bit AES key is generated.
    \item The plaintext is encrypted using AES-256-GCM, providing authenticated encryption (AEAD).
    \item The AES key (32~bytes) is encrypted using RSA with PKCS\#1 v1.5 padding~\cite{pkcs1_2016}, producing a wrapped key.
    \item The wrapped key, IV, ciphertext, and GCM tag are packaged into an envelope.
    \item Decryption reverses the process: RSA unwraps the AES key, which decrypts the ciphertext.
\end{enumerate}
This limits RSA to encrypting only 32~bytes regardless of data size, making throughput dominated by AES-GCM's hardware-accelerated performance.

\section{Experimental Setup}
\label{sec:setup}

\subsection{Test Environment}

All experiments were conducted on the following configuration:
\begin{itemize}
    \item \textbf{Processor:} Apple Silicon (ARM64), 10 performance cores
    \item \textbf{Memory:} 16~GB unified memory
    \item \textbf{Operating System:} macOS 26.4.1 (Darwin 25.4.0)
    \item \textbf{Python Version:} 3.14.4
    \item \textbf{Cryptography Library:} cryptography 46.0.7~\cite{cryptography_io}
    \item \textbf{Memory Profiler:} tracemalloc (Python standard library)~\cite{python_tracemalloc}
    \item \textbf{Concurrency:} threading + concurrent.futures~\cite{python_threading}
\end{itemize}

The Python cryptography library was used for AES-256-GCM operations in both implementations. All RSA operations (key generation, encryption, decryption, modular exponentiation) were implemented from scratch without relying on the library.

\subsection{Test Suite Design}

Six test suites were designed to evaluate both implementations:

\textbf{Test~1---Correctness Verification} (40~assertions): Validates mathematical correctness of both implementations. Tests cover modular exponentiation against Python's built-in \texttt{pow()}, RSA encrypt/decrypt roundtrips, byte-level operations with PKCS\#1 padding, cross-implementation interoperability, and hybrid RSA+AES-GCM roundtrips for payloads up to 5120~bytes.

\textbf{Test~2---Speed Benchmark} (50~samples): Measures wall-clock time using \texttt{time.perf\_counter()} for modular exponentiation at 1024/2048-bit, key generation, encrypt/decrypt operations, and cold-start versus cached key reuse.

\textbf{Test~3---Memory Usage} (20~samples): Uses \texttt{tracemalloc}~\cite{python_tracemalloc} to measure peak memory allocation during key generation, encryption, decryption, and batch operations.

\textbf{Test~4---Hybrid vs.\ Pure RSA}: Compares hybrid RSA+AES-GCM versus pure RSA chunked encryption at 1~KB, 10~KB, and 50~KB data sizes.

\textbf{Test~5---Key Size Scaling}: Measures performance across 512-bit, 1024-bit, and 2048-bit key sizes and computes growth factors.

\textbf{Test~6---Concurrency and Throughput}: Simulates multi-user scenarios using \texttt{ThreadPoolExecutor} with 1--10 threads, measuring throughput, latency, and error rates.

\subsection{Data Collection Methodology}

Timing uses \texttt{time.perf\_counter()} (nanosecond precision). Memory uses \texttt{tracemalloc.start()} and \texttt{tracemalloc.get\_traced\_memory()}. Each benchmark runs 50~samples (speed) or 20~samples (memory) and reports mean, median, standard deviation, minimum, and maximum. The test runner executes each test as a separate subprocess for clean memory state.

\section{Results and Analysis}
\label{sec:results}

\subsection{Correctness Verification}

Both implementations passed all 40 correctness assertions. Modular exponentiation results matched Python's built-in \texttt{pow(base, exp, mod)} for all test cases including 1024-bit operands. The optimized implementation's CRT-based decryption with fault verification produced identical results to the standard implementation. Cross-implementation tests (encrypting with standard RSA, decrypting with optimized RSA) also passed. Hybrid RSA+AES-GCM roundtrips succeeded for all payload sizes from 16~bytes to 5120~bytes.

\subsection{Key Generation Performance}

Key generation showed the most dramatic improvement in the optimized implementation, as summarized in Table~\ref{tab:keygen}.

\begin{table}[htbp]
\caption{Key Generation Time Comparison}
\begin{center}
\begin{tabular}{|c|c|c|c|}
\hline
\textbf{Key Size} & \textbf{Standard (ms)} & \textbf{Optimized (ms)} & \textbf{Speedup} \\
\hline
512-bit & 24.63 & 7.51 & 3.28$\times$ \\
\hline
1024-bit & 159.01 & 71.41 & 2.23$\times$ \\
\hline
2048-bit & 1644.70 & 462.88 & 3.55$\times$ \\
\hline
\end{tabular}
\label{tab:keygen}
\end{center}
\end{table}

The optimized implementation achieves a consistent 2.2--3.6$\times$ speedup, primarily attributed to small-primes pre-filtering in the optimized Miller-Rabin test. At 2048-bit key sizes, the optimized implementation generates keys in approximately 463~ms compared to over 1.6~seconds for the standard implementation.

\subsection{Encryption and Decryption Speed}

Table~\ref{tab:encdec} presents the encryption and decryption performance results.

\begin{table}[htbp]
\caption{Encryption and Decryption Time Comparison}
\begin{center}
\begin{tabular}{|c|c|c|c|c|}
\hline
\textbf{Operation} & \textbf{Key Size} & \textbf{Std (ms)} & \textbf{Opt (ms)} & \textbf{Speedup} \\
\hline
\multirow{3}{*}{Encrypt} & 512-bit & 0.015 & 0.118 & 0.13$\times$ \\
\cline{2-5}
 & 1024-bit & 0.034 & 0.305 & 0.11$\times$ \\
\cline{2-5}
 & 2048-bit & 0.116 & 1.045 & 0.11$\times$ \\
\hline
\multirow{3}{*}{Decrypt} & 512-bit & 0.727 & 0.730 & 1.00$\times$ \\
\cline{2-5}
 & 1024-bit & 3.401 & 4.154 & 0.82$\times$ \\
\cline{2-5}
 & 2048-bit & 23.096 & 18.703 & 1.23$\times$ \\
\hline
\end{tabular}
\label{tab:encdec}
\end{center}
\end{table}

For encryption, the standard implementation's naive square-and-multiply using Python's built-in arbitrary-precision integers is faster than the optimized implementation's Montgomery + sliding-window approach. This is because Python's \texttt{int} type uses highly optimized C-level arithmetic, and the overhead of Montgomery form conversion outweighs the theoretical benefit for the relatively small public exponent $e = 65537$ (only 17~bits).

For decryption, the CRT advantage becomes visible at larger key sizes. At 2048-bit, the optimized implementation is 1.23$\times$ faster because CRT splits the 2048-bit decryption into two 1024-bit operations.

The modular exponentiation benchmark in isolation (Table~\ref{tab:modexp}) shows this pattern clearly:

\begin{table}[htbp]
\caption{Modular Exponentiation Time (Isolated Benchmark)}
\begin{center}
\begin{tabular}{|c|c|c|c|}
\hline
\textbf{Operand Size} & \textbf{Standard (ms)} & \textbf{Optimized (ms)} & \textbf{Ratio} \\
\hline
512-bit & 0.733 & 2.034 & 0.36$\times$ \\
\hline
1024-bit & 3.257 & 9.180 & 0.35$\times$ \\
\hline
2048-bit & 23.352 & 69.548 & 0.34$\times$ \\
\hline
\end{tabular}
\label{tab:modexp}
\end{center}
\end{table}

The Montgomery multiplication implementation in Python is approximately 3$\times$ slower than Python's native \texttt{pow()} because per-operation Python overhead (object creation, function calls, dictionary lookups) dominates the actual arithmetic.

\subsection{Cold Start vs.\ Cached Key Performance}

The most impactful optimization in practical scenarios is key caching. With cold start (key generation + encrypt), the standard implementation takes approximately 178~ms per operation. With the optimized implementation's \texttt{KeyCache}, subsequent operations take only 0.30~ms---a speedup of 93$\times$. Under concurrent load with 100~sessions, the cached approach completes in 0.51~s versus 21.7~s for uncached, a 43$\times$ improvement.

\subsection{Hybrid RSA+AES-GCM vs.\ Pure RSA}

Table~\ref{tab:hybrid} confirms the overwhelming superiority of the hybrid approach.

\begin{table}[htbp]
\caption{Hybrid RSA+AES-GCM vs.\ Pure RSA Encryption Time}
\begin{center}
\begin{tabular}{|c|c|c|c|}
\hline
\textbf{Data Size} & \textbf{Pure RSA (ms)} & \textbf{Hybrid (ms)} & \textbf{Speedup} \\
\hline
1~KB & 2.632 & 0.318 & 8.3$\times$ \\
\hline
10~KB & 25.477 & 0.305 & 83.4$\times$ \\
\hline
50~KB & 126.702 & 0.318 & 398.0$\times$ \\
\hline
\end{tabular}
\label{tab:hybrid}
\end{center}
\end{table}

Pure RSA encrypts data in 117-byte chunks, requiring 9~RSA operations for 1~KB, 87 for 10~KB, and 435 for 50~KB. The hybrid approach performs exactly one RSA operation regardless of data size, with AES-GCM handling bulk encryption at near-constant time.

\subsection{Memory Usage}

Table~\ref{tab:memory} presents memory allocation patterns.

\begin{table}[htbp]
\caption{Memory Usage Comparison (2048-bit unless noted)}
\begin{center}
\begin{tabular}{|c|c|c|c|}
\hline
\textbf{Operation} & \textbf{Standard (KB)} & \textbf{Optimized (KB)} & \textbf{Change} \\
\hline
Key gen (2048-bit) & 4.30 & 4.30 & 0.0\% \\
\hline
Encrypt (1024-bit) & 1.39 & 3.22 & +131.9\% \\
\hline
Decrypt (1024-bit) & 1.40 & 2.99 & +113.5\% \\
\hline
mod\_exp (2048-bit) & 2.59 & 9.17 & +253.6\% \\
\hline
\end{tabular}
\label{tab:memory}
\end{center}
\end{table}

The optimized implementation allocates approximately 2--3$\times$ more memory per encryption/decryption operation due to Montgomery context storage and sliding-window pre-computation. However, in absolute terms, the overhead is small (5--9~KB per operation) and not a concern for modern server environments.

\subsection{Key Size Scaling}

Table~\ref{tab:scaling} shows how performance degrades with increasing key sizes.

\begin{table}[htbp]
\caption{Key Size Growth Factors (Standard Implementation)}
\begin{center}
\begin{tabular}{|c|c|c|}
\hline
\textbf{Transition} & \textbf{Encrypt Growth} & \textbf{Decrypt Growth} \\
\hline
512 $\rightarrow$ 1024 & 2.53$\times$ & 5.91$\times$ \\
\hline
1024 $\rightarrow$ 2048 & 3.61$\times$ & 4.42$\times$ \\
\hline
\end{tabular}
\label{tab:scaling}
\end{center}
\end{table}

Encryption time grows sub-quadratically because the public exponent $e = 65537$ remains constant. Decryption time grows approximately 4--5$\times$ per doubling, consistent with $O(n^2)$ complexity of modular exponentiation.

\subsection{Concurrency and Throughput}

Table~\ref{tab:concurrency} presents concurrent session simulation results.

\begin{table}[htbp]
\caption{Concurrent Session Throughput}
\begin{center}
\begin{tabular}{|c|c|c|c|c|}
\hline
\textbf{Threads} & \textbf{Std (sess/s)} & \textbf{Opt (sess/s)} & \textbf{Std Lat. (ms)} & \textbf{Opt Lat. (ms)} \\
\hline
1 & 291.5 & 226.3 & 3.42 & 4.41 \\
\hline
2 & 303.9 & 220.2 & 6.16 & 8.52 \\
\hline
5 & 303.1 & 219.2 & 10.83 & 16.22 \\
\hline
10 & 280.0 & 218.0 & 6.42 & 4.57 \\
\hline
\end{tabular}
\label{tab:concurrency}
\end{center}
\end{table}

Both implementations exhibit linear throughput scaling with no significant degradation at higher concurrency. The high-volume test (200~sessions, 10~threads) confirmed reliability with zero errors for both implementations. The cached versus uncached comparison under concurrency showed 0.51~s (cached) versus 21.7~s (uncached) for 100~sessions---a 43$\times$ improvement.

\section{Discussion}
\label{sec:discussion}

\subsection{Optimization Trade-offs}

The experimental results reveal that RSA optimization is not a simple matter of applying every available technique. Montgomery multiplication and sliding-window exponentiation, while theoretically superior, introduce Python-level overhead that makes them slower than Python's native C-implemented integer arithmetic for moderate key sizes. The CRT optimization, however, provides genuine speedup for decryption at 2048-bit and larger key sizes because it reduces the fundamental operation count.

The most impactful practical optimization is key caching. Our results show 43--93$\times$ speedup from key reuse, dwarfing all algorithmic improvements. This suggests that deployment-level optimizations (key management infrastructure, session key caching, connection pooling) may yield more practical benefit than algorithm-level improvements in interpreted language environments.

\subsection{Hybrid Encryption as a Necessity}

The 8.3--398$\times$ speedup from hybrid RSA+AES-GCM over pure RSA for bulk data confirms that hybrid encryption is not optional---it is essential. The hybrid approach's near-constant encryption time regardless of data size (0.31--0.33~ms for 1--50~KB) demonstrates that the RSA component becomes a fixed overhead for key wrapping, while AES-GCM handles bulk data at imperceptible speeds.

\subsection{Implications for Web Application Security}

These findings have direct implications for web application architects:
\begin{enumerate}
    \item RSA key pairs should be generated during server startup or through a background key rotation process, never on the request path.
    \item Hybrid encryption should be the default for any data exceeding 256~bytes.
    \item RSA key sizes should be chosen based on the threat model: 2048-bit for current requirements, with a migration plan to 3072-bit or PQC alternatives for long-term protection.
    \item The Python-specific finding that native arithmetic outperforms manual Montgomery multiplication suggests that high-performance RSA implementations should use compiled languages or hardware acceleration.
\end{enumerate}

\subsection{Limitations}

This study has several limitations. First, the implementations are written in Python, an interpreted language where per-operation overhead is high. Results would differ significantly in compiled languages. Second, the concurrency tests use Python's GIL-affected threading, limiting true parallelism. Third, only three key sizes (512, 1024, 2048) were tested due to the time required for from-scratch key generation. Fourth, the study does not evaluate side-channel resistance beyond fault-attack verification.

\section{Conclusion}
\label{sec:conclusion}

This study presented a comprehensive performance evaluation of the RSA cryptosystem in web application contexts, comparing a standard baseline implementation against an optimized variant incorporating CRT-based decryption, sliding-window modular exponentiation, Montgomery multiplication, and key caching. Six benchmark test suites covering correctness, speed, memory, hybrid encryption, key size scaling, and concurrent throughput were designed and executed.

The key findings are: (1)~Key caching provides the largest practical speedup (43--93$\times$), validating the importance of key management infrastructure. (2)~CRT-based decryption provides a measurable but modest improvement (1.23$\times$ at 2048-bit). (3)~Montgomery multiplication and sliding-window exponentiation do not improve performance in Python due to interpreter overhead. (4)~Hybrid RSA+AES-GCM encryption is 8.3--398$\times$ faster than pure RSA for bulk data. (5)~RSA performance scales approximately quadratically with key size. (6)~Both implementations handle concurrent load reliably with zero errors.

These results support three recommendations for practitioners: adopt hybrid encryption for all data exceeding a few hundred bytes, implement key caching to eliminate the dominant key-generation bottleneck, and plan for post-quantum migration as RSA's security margin narrows. Future work should replicate these benchmarks in compiled languages, evaluate post-quantum alternatives (CRYSTALS-Kyber~\cite{kyber2023}, CRYSTALS-Dilithium~\cite{dilithium2022}) alongside RSA, and measure energy consumption for mobile and edge deployments.

\bibliographystyle{IEEEtran}
\begin{thebibliography}{20}

\bibitem{rivest1978}
R.~L. Rivest, A.~Shamir, and L.~Adleman, ``A method for obtaining digital signatures and public-key cryptosystems,'' \textit{Communications of the ACM}, vol.~21, no.~2, pp.~120--126, 1978.

\bibitem{boneh1999}
D.~Boneh, ``Twenty years of attacks on the RSA cryptosystem,'' \textit{Notices of the American Mathematical Society}, vol.~46, no.~2, pp.~203--213, 1999.

\bibitem{coppersmith1997}
D.~Coppersmith, ``Small solutions to polynomial equations, and low exponent RSA vulnerabilities,'' \textit{Journal of Cryptology}, vol.~10, no.~4, pp.~233--260, 1997.

\bibitem{shor1997}
P.~W. Shor, ``Polynomial-time algorithms for prime factorization and discrete logarithms on a quantum computer,'' \textit{SIAM Journal on Computing}, vol.~26, no.~5, pp.~1484--1509, 1997.

\bibitem{wiener1990}
M.~Wiener, ``Cryptanalysis of short RSA secret exponents,'' \textit{IEEE Trans. Inf. Theory}, vol.~36, no.~3, pp.~553--558, 1990.

\bibitem{boneh2001}
D.~Boneh, R.~A. DeMillo, and R.~J. Lipton, ``On the importance of eliminating errors in cryptographic computations,'' \textit{Journal of Cryptology}, vol.~14, no.~2, pp.~101--119, 2001.

\bibitem{garner1959}
H.~L. Garner, ``The residue number system,'' \textit{IRE Trans. Electronic Computers}, vol.~EC-8, no.~2, pp.~140--147, 1959.

\bibitem{montgomery1985}
P.~L. Montgomery, ``Modular multiplication without trial division,'' \textit{Mathematics of Computation}, vol.~44, no.~170, pp.~519--521, 1985.

\bibitem{pkcs1_2016}
RSA Laboratories, ``PKCS \#1: RSA cryptography specifications version 2.2,'' RFC 8017, IETF, 2016.

\bibitem{nist2024}
National Institute of Standards and Technology, ``Post-quantum cryptography standardization,'' NIST, 2024. [Online]. Available: \url{https://csrc.nist.gov/projects/post-quantum-cryptography}

\bibitem{kyber2023}
R.~Avanzi \textit{et al.}, ``CRYSTALS-Kyber: A CCA-secure module-lattice-based KEM,'' in \textit{Proc. IEEE Symp. Security and Privacy}, 2023.

\bibitem{dilithium2022}
S.~Bai \textit{et al.}, ``CRYSTALS-Dilithium: A lattice-based digital signature scheme,'' \textit{IACR Trans. Cryptographic Hardware and Embedded Systems}, vol.~2022, no.~1, pp.~218--244, 2022.

\bibitem{bellare1995}
M.~Bellare and P.~Rogaway, ``Optimal asymmetric encryption padding,'' in \textit{Advances in Cryptology --- EUROCRYPT '94}, LNCS, vol.~950, Springer, pp.~92--111, 1995.

\bibitem{bernstein2010}
D.~J. Bernstein, ``Grover vs.\ McEliece,'' in \textit{Post-Quantum Cryptography}, LNCS, vol.~6061, Springer, pp.~89--98, 2010.

\bibitem{menezes2001}
A.~J. Menezes, P.~C. van Oorschot, and S.~A. Vanstone, \textit{Handbook of Applied Cryptography}. CRC Press, 5th~ed., 2001.

\bibitem{shoup2008}
V.~Shoup, \textit{A Computational Introduction to Number Theory and Algebra}, Version~2.3.1, Cambridge Univ.\ Press, 2008. [Online]. Available: \url{https://shoup.net/ntb/}

\bibitem{python_tracemalloc}
Python Software Foundation, ``The Python standard library --- tracemalloc: Trace memory allocations,'' Python~3 Documentation, 2025. [Online]. Available: \url{https://docs.python.org/3/library/tracemalloc.html}

\bibitem{python_threading}
Python Software Foundation, ``The Python standard library --- threading: Thread-based parallelism,'' Python~3 Documentation, 2025. [Online]. Available: \url{https://docs.python.org/3/library/threading.html}

\bibitem{cryptography_io}
Cryptography Software Foundation, ``cryptography --- A Python library exposing cryptographic recipes and primitives,'' Version~46.0.7, 2025. [Online]. Available: \url{https://cryptography.io/}

\bibitem{kelsey1999}
J.~A. Kelsey, B.~Schneier, and N.~Ferguson, ``Yarrow-160: Notes on the design and analysis of the Yarrow cryptographic pseudorandom number generator,'' in \textit{Proc. Sixth Annual Workshop on Selected Areas in Cryptography}, 1999.

\end{thebibliography}

\end{document}
